% ==================================================================
% id:       cond-002
% topic:    הסתברות מותנית ואי-תלות
% subtopic: נוסחת ההסתברות השלמה ונוסחת בייס
% solution: full
% ==================================================================
בשעשוע טלוויזיה יש שלוש דלתות, ומאחורי אחת מהן (באקראי) מכונית. המתמודד בוחר דלת. המנחה, שיודע איפה המכונית, פותח אחת משתי הדלתות האחרות שאין מאחוריה מכונית (אם יש לו שתי אפשרויות, הוא בוחר ביניהן באקראי). עכשיו המתמודד יכול להחליף את בחירתו לדלת הסגורה השנייה. האם כדאי לו להחליף?
\begin{solbox}
כן: אם המתמודד מחליף, הוא זוכה בהסתברות $\frac23$, ואם הוא לא מחליף, בהסתברות $\frac13$.

בגלל הסימטריה, מספיק לבדוק את המקרה שבו המתמודד בחר בדלת 1 והמנחה פתח את דלת 3. נסמן $C_i$ = „המכונית מאחורי דלת $i$” ($i=1,2,3$), ו-$H$ = „המנחה פתח את דלת 3”. המכונית מוצבת באקראי, ולכן $\PP(C_1)=\PP(C_2)=\PP(C_3)=\frac13$, והמאורעות $C_1,C_2,C_3$ הם חלוקה של $\Omega$. לפי הכללים של המנחה:
\[
\PP(H\mid C_1)=\tfrac12,\qquad \PP(H\mid C_2)=1,\qquad \PP(H\mid C_3)=0.
\]
(אם המכונית מאחורי דלת 1, המנחה בוחר באקראי בין דלתות 2 ו-3; אם היא מאחורי דלת 2, הוא חייב לפתוח את דלת 3; ואת הדלת שמאחוריה המכונית הוא לא פותח.)

לפי נוסחת ההסתברות השלמה
\[
\PP(H)=\tfrac12\cdot\tfrac13+1\cdot\tfrac13+0\cdot\tfrac13=\tfrac12.
\]
לפי נוסחת בייס
\[
\PP(C_1\mid H)=\frac{\PP(H\mid C_1)\PP(C_1)}{\PP(H)}=\frac{\frac12\cdot\frac13}{\frac12}=\frac13,\qquad
\PP(C_2\mid H)=\frac{1\cdot\frac13}{\frac12}=\ans{\frac23}.
\]
כלומר, אחרי שהמנחה פתח את דלת 3, המכונית מאחורי הדלת שהמתמודד בחר בהסתברות $\frac13$ בלבד, ומאחורי הדלת השנייה בהסתברות $\frac23$. לכן כדאי להחליף.
\end{solbox}
